\documentclass[conference]{IEEEtran}
\IEEEoverridecommandlockouts
% The preceding line is only needed to identify funding in the first footnote. If that is unneeded, please comment it out.
\usepackage{cite}
\usepackage{amsmath,amssymb,amsfonts}
\usepackage{algorithmic}
\usepackage{graphicx}
\usepackage{textcomp}
\usepackage{xcolor}
\def\BibTeX{{\rm B\kern-.05em{\sc i\kern-.025em b}\kern-.08em
    T\kern-.1667em\lower.7ex\hbox{E}\kern-.125emX}}
\begin{document}

\title{
\thanks{Identify applicable funding agency here. If none, delete this.}
}

\author{\IEEEauthorblockN{1\textsuperscript{st} Given Name Surname}
\IEEEauthorblockA{\textit{dept. name of organization (of Aff.)} \\
\textit{name of organization (of Aff.)}\\
City, Country \\
email address or ORCID}
\and
\IEEEauthorblockN{2\textsuperscript{nd} Given Name Surname}
\IEEEauthorblockA{\textit{dept. name of organization (of Aff.)} \\
\textit{name of organization (of Aff.)}\\
City, Country \\
email address or ORCID}}

\maketitle

\begin{abstract}
This survey provides a comprehensive overview of the current state of research in multi-camera multi-object tracking (MOT), a field that has experienced significant growth in recent years due to advancements in computer vision, machine learning, and increasing demand for intelligent surveillance systems. The paper reviews the fundamental concepts, traditional approaches, and recent developments in MOT, including deep learning-based methods, transformer-based architectures, and efficient real-time algorithms. The scope of this survey encompasses various aspects of MOT, ranging from background and fundamentals to evaluation metrics and methodologies, as well as applications in automated driving systems.

The key findings of this survey highlight the importance of benchmarking efforts, such as the MOTChallenge, in advancing the field of MOT. Additionally, the paper discusses open challenges and future directions in MOT research, including the need for standardized benchmarks, improved detection-aware and polygon-based tracking methods, and more efficient memory-augmented and online techniques. The contributions of this survey include a thorough analysis of relevant papers, identification of common themes and technical details, and a summary of key findings that can inform future research in MOT.

This comprehensive review aims to provide researchers and practitioners with a detailed understanding of the current state of MOT research, its applications, and future directions. By highlighting the scope and contributions of this survey, we hope to facilitate further advancements in the field and encourage continued innovation in MOT. Ultimately, this survey seeks to contribute to the development of more effective and efficient tracking methods that can be applied in a variety of real-world scenarios, from security and surveillance to autonomous vehicles and robotics.
\end{abstract}

\begin{IEEEkeywords}
large language models, multimodal learning, natural language processing, machine learning, artificial intelligence
\end{IEEEkeywords}


\section{Introduction to Multi-Object Tracking}
The field of multi-object tracking has experienced significant growth in recent years, driven by advancements in computer vision, machine learning, and the increasing demand for intelligent surveillance systems \cite{MOTChallenge}. At its core, multi-object tracking involves the detection and tracking of multiple objects across a sequence of images or video frames, with applications ranging from security and surveillance to autonomous vehicles and robotics. The complexity of this task is further exacerbated when dealing with multiple cameras, where the challenges of non-overlapping cameras and inter-camera object tracking must be addressed \cite{Challenge of Multi-Camera Tracking}. 

A key theme that emerges from recent research is the importance of standardized benchmarks for evaluating multi-object tracking methods. Frameworks such as MOTChallenge \cite{MOTChallenge} and Tracking the Trackers \cite{Tracking the Trackers} have been instrumental in providing a common platform for comparing and evaluating different approaches, highlighting the need for rigorous evaluation protocols to advance the field. The introduction of novel end-to-end approaches, such as MCTR \cite{MCTR}, which leverages transformers for multi-camera tracking, demonstrates the potential of deep learning-based methods in tackling this complex task.

The use of graph-based models, as proposed by An equalised global graphical model-based approach \cite{An equalised global graphical model-based approach}, offers an alternative perspective on inter-camera object tracking, emphasizing the importance of optimizing similarity metrics for effective tracking across non-overlapping cameras. This diversity in approaches underscores the vibrant nature of current research, with different methodologies being explored to address the unique challenges posed by multi-camera multi-object tracking.

Despite these advances, significant challenges remain, including the lack of large-scale datasets and annotations, which hinders the development and evaluation of more sophisticated tracking algorithms \cite{MOTChallenge}. Furthermore, the technical complexities associated with non-overlapping cameras and inter-camera object tracking continue to pose substantial hurdles \cite{Challenge of Multi-Camera Tracking}. Addressing these challenges will require continued innovation in both algorithmic design and the development of standardized evaluation frameworks.

The implications of advancements in multi-object tracking are far-reaching, with potential applications in enhancing public safety through intelligent surveillance systems, improving the efficiency of traffic management, and facilitating the development of autonomous vehicles \cite{MOTChallenge}. As research continues to evolve, it is crucial to consider not only the technical aspects of multi-object tracking but also its societal implications and ethical considerations. By doing so, we can ensure that these technologies are developed and deployed in a manner that benefits society as a whole.

In conclusion, the field of multi-object tracking, particularly in the context of multiple cameras, represents a dynamic and challenging area of research. Through the development of standardized benchmarks, the exploration of novel deep learning-based approaches, and the addressing of technical complexities, researchers are making significant strides towards overcoming the hurdles associated with this task. As we move forward, it is essential to maintain a holistic perspective, considering both the technical and societal dimensions of multi-object tracking, to fully realize its potential and contribute meaningfully to the broader community.

\section{Background and Fundamentals of MOT}
The Multi-Object Tracking (MOT) problem has garnered significant attention in recent years due to its academic and commercial potential, particularly with advancements in autonomous driving. A thorough analysis of relevant papers reveals that the MOTChallenge benchmark has been instrumental in pushing the performance of computer vision algorithms for multiple object tracking \cite{MOTChallenge: A Benchmark for Multi-Object Tracking}. Recent reviews have highlighted the importance of tackling challenges such as occlusion, appearance changes, and small object detection in MOT \cite{Multiple Object Tracking: A Review}, \cite{A Comprehensive Review of Multiple Object Tracking Techniques}. The use of attention mechanisms, graph convolutional neural networks, siamese networks, and motion prediction with LSTM has been explored to address these challenges \cite{Multiple Object Tracking: A Review}.

The need for standardized benchmarks to evaluate the performance of MOT algorithms is a recurring theme across various papers \cite{MOTChallenge: A Benchmark for Multi-Object Tracking}, \cite{MOT16: A Benchmark for Multi-Object Tracking}, \cite{MOT20: A Benchmark for Multi-Object Tracking in Crowded Scenes}. Benchmark datasets such as MOT15, MOT16, MOT17, and MOT20 have been introduced to evaluate the performance of multi-object trackers in various scenarios \cite{MOTChallenge: A Benchmark for Multi-Object Tracking}, \cite{MOT16: A Benchmark for Multi-Object Tracking}, \cite{MOT20: A Benchmark for Multi-Object Tracking in Crowded Scenes}. The importance of addressing challenges such as occlusion, appearance changes, and small object detection is highlighted in multiple papers \cite{Multiple Object Tracking: A Review}, \cite{A Comprehensive Review of Multiple Object Tracking Techniques}.

While some papers focus on the development of new benchmarks for evaluating MOT algorithms \cite{MOTChallenge: A Benchmark for Multi-Object Tracking}, \cite{MOT16: A Benchmark for Multi-Object Tracking}, \cite{MOT20: A Benchmark for Multi-Object Tracking in Crowded Scenes}, others provide comprehensive reviews of existing approaches and techniques \cite{Multiple Object Tracking: A Review}, \cite{A Comprehensive Review of Multiple Object Tracking Techniques}. The use of simple IOU matching-based CNN networks versus more complex approaches such as graph convolutional neural networks is a point of contrast between different papers \cite{Multiple Object Tracking: A Review}. The MOTChallenge benchmark uses a consistent annotation protocol to label objects in videos, providing a standardized framework for evaluating MOT algorithms \cite{MOTChallenge: A Benchmark for Multi-Object Tracking}.

The use of siamese networks and motion prediction with LSTM has been explored to address challenges such as appearance changes and small object detection \cite{Multiple Object Tracking: A Review}. Graph convolutional neural networks have been used to model the interrelation between tracklets in MOT \cite{Multiple Object Tracking: A Review}. Despite these advancements, the MOT problem remains challenging due to factors such as abrupt appearance changes, severe object occlusions, and small object detection difficulty \cite{Multiple Object Tracking: A Review}, \cite{A Comprehensive Review of Multiple Object Tracking Techniques}. The need for more accurate and efficient algorithms that can handle crowded scenes and complex scenarios is a recurring challenge highlighted in multiple papers \cite{MOTChallenge: A Benchmark for Multi-Object Tracking}, \cite{MOT16: A Benchmark for Multi-Object Tracking}, \cite{MOT20: A Benchmark for Multi-Object Tracking in Crowded Scenes}.

The evaluation of MOT algorithms using standardized benchmarks is limited by the availability of high-quality annotated data \cite{MOTChallenge: A Benchmark for Multi-Object Tracking}. In conclusion, the analysis of relevant papers reveals that the development of new benchmarks and the use of deep learning-based approaches are common themes in MOT research. However, limitations and challenges remain, including the need for more accurate and efficient algorithms that can handle crowded scenes and complex scenarios. As the field continues to evolve, it is essential to address these challenges and develop more robust and effective MOT algorithms \cite{Multiple Object Tracking: A Review}, \cite{A Comprehensive Review of Multiple Object Tracking Techniques}.

\section{Traditional Approaches to Multi-Object Tracking}
Traditional approaches to multi-object tracking have been extensively studied in the literature, with various papers contributing to the understanding of this complex problem. The paper "Challenge of Multi-Camera Tracking" \cite{Challenge of Multi-Camera Tracking} highlights the differences between single-camera and multi-camera tracking, emphasizing the need for new technologies and system architectures to address the challenges of multi-camera tracking. Another notable work is "Tracking the Trackers: An Analysis of the State of the Art in Multiple Object Tracking" \cite{Tracking the Trackers: An Analysis of the State of the Art in Multiple Object Tracking}, which presents a benchmark for evaluating multiple object tracking methods, providing an in-depth analysis of state-of-the-art trackers and identifying current trends and weaknesses.

Several papers have proposed various approaches to improve multi-camera object tracking performance. For instance, "An equalised global graphical model-based approach for multi-camera object tracking" \cite{An equalised global graphical model-based approach for multi-camera object tracking} proposes a global graph model with an improved similarity metric, while "MCTR: Multi Camera Tracking Transformer" \cite{MCTR: Multi Camera Tracking Transformer} introduces an end-to-end approach, leveraging detectors like DEtector TRansformer (DETR) to produce detections and detection embeddings for each camera view. A comprehensive literature review is provided in "Multiple Object Tracking: A Literature Review" \cite{Multiple Object Tracking: A Literature Review}, discussing key aspects, existing approaches, and future research directions.

A common theme that emerges from the analysis of these papers is the importance of multi-camera tracking in various applications. The challenges and limitations of traditional approaches to multi-object tracking are also highlighted, including abrupt appearance changes, severe object occlusions, and the need for new technologies and system architectures. Furthermore, the papers stress the importance of standardized evaluation and benchmarking to compare the performance of different tracking methods. Graphical models and end-to-end approaches are proposed as potential solutions to improve multi-camera object tracking performance.

The technical details and methodologies used in these papers include graphical models, detectors, and detection embeddings. For example, "An equalised global graphical model-based approach for multi-camera object tracking" \cite{An equalised global graphical model-based approach for multi-camera object tracking} uses a global graph model with an improved similarity metric, while "MCTR: Multi Camera Tracking Transformer" \cite{MCTR: Multi Camera Tracking Transformer} leverages detectors like DEtector TRansformer (DETR) to produce detections and detection embeddings. Benchmarking and evaluation metrics, such as precision, recall, and F1 score, are also used to evaluate the performance of tracking methods.

Despite the progress made in traditional approaches to multi-object tracking, several limitations and challenges remain. Abrupt appearance changes and severe object occlusions are still challenging to handle, and scalable and computationally efficient algorithms are needed to handle large-scale multi-camera tracking applications. Moreover, there is a need for standardized evaluation and benchmarking to compare the performance of different tracking methods. These findings, themes, and challenges provide a foundation for understanding traditional approaches to multi-object tracking and can inform the development of more effective and efficient tracking algorithms.

The implications of these developments are significant, as they have the potential to improve the accuracy and efficiency of multi-object tracking systems in various applications, such as surveillance, robotics, and autonomous vehicles. Furthermore, the connections between traditional approaches and more recent deep learning-based methods are becoming increasingly important, as researchers seek to combine the strengths of both paradigms to develop more effective and robust tracking algorithms. Overall, the study of traditional approaches to multi-object tracking remains a vital area of research, with ongoing developments and innovations continuing to shape the field.

\section{Deep Learning-Based Methods for MOT}
Deep learning-based methods for multi-object tracking (MOT) have gained significant attention in recent years, with numerous papers presenting innovative approaches to address the challenges of tracking multiple objects across frames in a video sequence. A key finding in this area is the development of end-to-end tracking methods, such as those presented in \cite{End-to-End Multi-Object Tracking with Global Response Map}, which take an image sequence as input and output the located and tracked objects directly. This approach has been made possible by the use of deep learning-based architectures, including convolutional neural networks (CNNs), recurrent neural networks (RNNs), and graph convolutional neural networks (GCNs), as discussed in \cite{Deep Learning in Video Multi-Object Tracking: A Survey} and \cite{Multiple Object Tracking in Recent Times: A Literature Review}.

The tracking-by-detection paradigm is a common theme in MOT methods, where objects are first detected and then associated across frames. Appearance features, such as those extracted from CNNs, play a crucial role in cross-frame association, allowing for the identification of objects across different frames. The importance of benchmarking is also highlighted in several papers, with the MOTChallenge benchmark \cite{MOTChallenge: A Benchmark for Single-Camera Multiple Target Tracking} being widely adopted in the field to evaluate the performance of MOT methods.

However, there are contrasting viewpoints or approaches to deep learning-based MOT, including online vs. offline tracking and single-camera vs. multi-camera tracking. Online tracking methods, such as those presented in \cite{End-to-End Multi-Object Tracking with Global Response Map}, process the video stream in real-time, while offline tracking methods, such as those discussed in \cite{Deep Learning in Video Multi-Object Tracking: A Survey}, process the entire video sequence at once. Single-camera MOT methods are also prevalent, but multi-camera MOT methods, such as those presented in \cite{Multiple Object Tracking in Recent Times: A Literature Review}, offer a more comprehensive approach to tracking objects across multiple cameras.

The technical details and methodologies used in deep learning-based MOT methods vary widely. Network architectures, such as CNNs, RNNs, and GCNs, are commonly used for feature extraction and object association. Loss functions, such as the intersection over union (IoU) loss \cite{End-to-End Multi-Object Tracking with Global Response Map}, are also crucial in training MOT models. Evaluation metrics, including the clear motion threshold (CMT), fragmentation (FRAG), and identity switch (IDS) metrics, are used to assess the performance of MOT methods.

Despite the significant progress made in deep learning-based MOT, there are still several limitations and challenges that need to be addressed. Occlusion handling remains a significant challenge, with many methods struggling to handle cases where objects are partially or fully occluded. Scalability is also an issue, with some papers highlighting the need for more scalable MOT methods that can handle large numbers of objects and high-resolution video sequences \cite{Multiple Object Tracking in Recent Times: A Literature Review}. The performance of deep learning-based MOT methods in real-world applications, such as self-driving cars or surveillance systems, remains a significant challenge due to the complexity and variability of real-world environments.

In conclusion, deep learning-based methods for MOT have made significant progress in recent years, with innovative approaches being developed to address the challenges of tracking multiple objects across frames. However, there are still several limitations and challenges that need to be addressed, including occlusion handling, scalability, and performance in real-world applications. Further research is needed to develop more robust and scalable MOT methods that can handle the complexities of real-world environments.

\section{Transformer-Based Architectures for MOT}
Transformer-Based Architectures for MOT have recently gained significant attention in the field of Multi-Object Tracking, with several notable contributions and advancements being reported in the literature. The proposal of TransTrack, a simple yet efficient scheme for solving MOT problems using a transformer architecture, has been particularly influential \cite{TransTrack}. Similarly, the development of MotionTrack, an end-to-end transformer-based MOT algorithm that utilizes LiDAR-camera fusion to track objects with multiple classes, has demonstrated the potential of transformer-based architectures in achieving state-of-the-art performance \cite{MotionTrack}.

A key theme that emerges from the analysis of these papers is the importance of attention mechanisms in transformer-based MOT algorithms. The use of modified attention mechanisms, such as the query-key mechanism used in TransTrack \cite{TransTrack}, has been shown to be effective in improving tracking performance. Furthermore, the introduction of a modified attention mechanism for data association in MOT has highlighted the need for ongoing research and experimentation in this area \cite{ModifiedAttention}. The achievement of state-of-the-art MOTA scores using transformer-based architectures, such as 73.20\% on the MOT17 benchmark \cite{MOT17}, demonstrates the competitiveness of these approaches in MOT.

The use of transformer architectures for MOT is a growing trend, with many researchers exploring their potential for improving tracking performance. The need for efficient and scalable architectures, such as the use of butterfly transform operations to reduce computational complexity \cite{ButterflyTransform}, is also a common theme. Additionally, the increasing use of multi-modality sensor inputs, such as LiDAR-camera fusion, to improve tracking performance in autonomous driving scenarios \cite{LiDARCamerFusion} highlights the importance of leveraging multiple sources of information to enhance MOT performance.

While there are many promising results and contributions in this area, there are also several limitations and challenges that need to be addressed. The computational complexity of transformer-based architectures can be a challenge for real-time tracking applications \cite{ComputationalComplexity}. Moreover, the scalability of these algorithms is an ongoing research topic, with many papers proposing techniques to reduce computational complexity \cite{Scalability}. The problem of data association remains a challenging one, particularly in scenarios with multiple objects and occlusions \cite{DataAssociation}. Finally, the availability and quality of benchmark datasets can be a limitation, particularly for evaluating the performance of MOT algorithms in real-world scenarios \cite{BenchmarkDatasets}.

The implications of these developments are significant, as they highlight the potential of transformer-based architectures to improve the performance and scalability of MOT algorithms. The connections between these papers demonstrate a growing trend towards the use of transformer-based architectures in MOT, with many researchers exploring their potential for improving tracking performance. Overall, the analysis of these papers provides a comprehensive overview of the current state-of-the-art in transformer-based architectures for MOT, highlighting key themes, developments, and implications for future research in this area. As the field continues to evolve, it is likely that we will see further advancements in the use of transformer-based architectures for MOT, with potential applications in autonomous driving, surveillance, and other areas where multi-object tracking is critical \cite{FutureDirections}.

\section{Efficient and Real-Time MOT Algorithms}
Efficient and Real-Time MOT Algorithms have been a crucial focus of research in the field of Multi Camera Multi Object Tracking, with numerous papers contributing to the development of robust and accurate tracking systems. A comprehensive review of recent advances in this area highlights the challenges and opportunities present in MOT, including the introduction of new trackers such as BoT-SORT \cite{BoT-SORT paper}, which combines motion and appearance information with camera-motion compensation and a more accurate Kalman filter state vector to improve tracking accuracy. Furthermore, benchmarks like MOTChallenge \cite{MOTChallenge paper} have been proposed to provide a standardized framework for evaluating MOT methods, emphasizing the importance of addressing challenges such as occlusion, similar appearance, small object detection difficulty, ID switching, and motion prediction.

The surveyed papers emphasize the significance of deep learning models in MOT algorithms, particularly in the four main steps: object detection, feature extraction, tracking, and association \cite{Deep Learning for MOT}. The use of convolutional neural networks (CNNs) and recurrent neural networks (RNNs) has become a common theme in recent MOT research, with many papers proposing solutions that balance accuracy and computational efficiency. For instance, the employment of attention mechanisms \cite{Attention Mechanisms for MOT} and graph convolutional neural networks \cite{Graph Convolutional Neural Networks for MOT} have been explored to improve tracking performance.

While some papers focus on the development of new tracking algorithms, others stress the importance of benchmarking and evaluation. The need for efficient and real-time MOT algorithms is a recurring challenge, with many current solutions requiring significant computational resources. Technical details such as camera-motion compensation \cite{Camera-Motion Compensation} and the use of siamese networks \cite{Siamese Networks for MOT} have been proposed to address these challenges. Despite recent advances, MOT remains a challenging problem, particularly in scenarios with severe occlusion or similar appearance, highlighting the need for more robust and accurate tracking algorithms.

The limitations and challenges present in current MOT research include the requirement for larger and more diverse benchmark datasets to evaluate MOT methods and promote further research in this field \cite{Need for Larger Datasets}. Additionally, the development of more efficient and real-time tracking algorithms is an ongoing challenge. To address these challenges, future research directions may involve investigating emerging techniques such as attention mechanisms or graph convolutional neural networks to improve MOT performance. The creation of larger and more diverse benchmark datasets will also be essential in promoting further research and evaluating the effectiveness of proposed MOT methods.

In conclusion, recent advances in Efficient and Real-Time MOT Algorithms have highlighted the importance of addressing challenges such as occlusion and similar appearance, with deep learning models and other techniques being employed to improve tracking accuracy. However, limitations and challenges remain, emphasizing the need for more robust and accurate tracking algorithms, larger and more diverse benchmark datasets, and further research into emerging techniques to promote progress in this field \cite{Future Research Directions}. The connections between these developments and their implications for real-world applications will be crucial in shaping the future of MOT research.

\section{Memory-Augmented and Online MOT Techniques}
Memory-Augmented and Online MOT Techniques 
====================================================

The development of memory-augmented and online Multi-Object Tracking (MOT) techniques has been a significant focus of research in recent years. Key findings and contributions in this area include the exploration of attention mechanisms, graph convolutional neural networks, and siamese networks to tackle challenges in MOT \cite{Multiple Object Tracking in Recent Times: A Literature Review}, \cite{Multiple Object Tracking: A Literature Review}. Online MOT methods have achieved desirable tracking performance but often at the cost of slow tracking speeds \cite{Real-time Online Multi-Object Tracking in Compressed Domain}. Techniques such as heuristic search for structural constraints \cite{Heuristic Search for Structural Constraints in Data Association} and discriminative appearance modeling with multi-track pooling \cite{Discriminative Appearance Modeling with Multi-track Pooling for Real-time Multi-object Tracking} have been proposed to improve data association and tracking accuracy. Furthermore, the use of compressed domain tracking and key frame detection has been shown to achieve real-time online MOT with comparable performance to state-of-the-art methods \cite{Real-time Online Multi-Object Tracking in Compressed Domain}.

A common theme across these papers is the importance of appearance modeling and data association in MOT \cite{Multiple Object Tracking in Recent Times: A Literature Review}, \cite{Multiple Object Tracking: A Literature Review}, \cite{Heuristic Search for Structural Constraints in Data Association}, \cite{Discriminative Appearance Modeling with Multi-track Pooling for Real-time Multi-object Tracking}. The use of deep learning techniques, such as CNNs and LSTMs, is also prevalent in MOT research \cite{Multiple Object Tracking in Recent Times: A Literature Review}, \cite{Multiple Object Tracking: A Literature Review}, \cite{Real-time Online Multi-Object Tracking in Compressed Domain}, \cite{Discriminative Appearance Modeling with Multi-track Pooling for Real-time Multi-object Tracking}. Online MOT methods are increasingly being explored to address the need for real-time tracking \cite{Real-time Online Multi-Object Tracking in Compressed Domain}, \cite{Discriminative Appearance Modeling with Multi-track Pooling for Real-time Multi-object Tracking}. However, handling similar appearances and occlusions in crowded scenes remains a significant challenge \cite{Multiple Object Tracking in Recent Times: A Literature Review}, \cite{Multiple Object Tracking: A Literature Review}, \cite{Heuristic Search for Structural Constraints in Data Association}, \cite{Discriminative Appearance Modeling with Multi-track Pooling for Real-time Multi-object Tracking}.

Different papers have proposed contrasting viewpoints or approaches to address these challenges. For instance, some focus on improving data association using structural constraints \cite{Heuristic Search for Structural Constraints in Data Association}, while others propose novel appearance modeling techniques \cite{Discriminative Appearance Modeling with Multi-track Pooling for Real-time Multi-object Tracking}. The use of compressed domain tracking \cite{Real-time Online Multi-Object Tracking in Compressed Domain} is an alternative approach to traditional MOT methods. Additionally, papers differ in their evaluation metrics and benchmark datasets, with some using MOTA \cite{Heuristic Search for Structural Constraints in Data Association} and others using more comprehensive evaluations \cite{Multiple Object Tracking in Recent Times: A Literature Review}, \cite{Multiple Object Tracking: A Literature Review}. From a technical standpoint, heuristic search for structural constraints involves searching for a maximum match set under minimum structural cost \cite{Heuristic Search for Structural Constraints in Data Association}, while discriminative appearance modeling with multi-track pooling uses a novel pooling module to update the memory of each object in the scene \cite{Discriminative Appearance Modeling with Multi-track Pooling for Real-time Multi-object Tracking}. Compressed domain tracking involves dividing frames into key and non-key frames and using motion information to propagate objects \cite{Real-time Online Multi-Object Tracking in Compressed Domain}. Siamese networks and graph convolutional neural networks are used for appearance modeling and data association \cite{Multiple Object Tracking in Recent Times: A Literature Review}, \cite{Multiple Object Tracking: A Literature Review}.

Despite the advancements, MOT research is hindered by several limitations and challenges. Handling similar appearances, occlusions, and crowded scenes remains a significant challenge \cite{Multiple Object Tracking in Recent Times: A Literature Review}, \cite{Multiple Object Tracking: A Literature Review}, \cite{Heuristic Search for Structural Constraints in Data Association}, \cite{Discriminative Appearance Modeling with Multi-track Pooling for Real-time Multi-object Tracking}. Online MOT methods face the challenge of balancing tracking performance and speed \cite{Real-time Online Multi-Object Tracking in Compressed Domain}. Furthermore, the choice of evaluation metrics and benchmark datasets can significantly impact the comparison of different MOT techniques \cite{Multiple Object Tracking in Recent Times: A Literature Review}, \cite{Multiple Object Tracking: A Literature Review}, \cite{Heuristic Search for Structural Constraints in Data Association}, \cite{Discriminative Appearance Modeling with Multi-track Pooling for Real-time Multi-object Tracking}. The implications of these developments are significant, as they can be applied to various real-world applications such as surveillance, autonomous vehicles, and robotics. The connections between different techniques and approaches highlight the need for a comprehensive framework that integrates multiple strategies to achieve robust and efficient MOT. Overall, memory-augmented and online MOT techniques have shown promising results, but further research is needed to address the existing challenges and limitations in this field.

\section{Evaluation Metrics and Methodologies for MOT}
The evaluation metrics and methodologies for Multi-Object Tracking (MOT) have undergone significant development in recent years, driven by the need for standardized benchmarks and objective measures of performance \cite{MOTChallenge: A Benchmark for Multi-Object Tracking}. The establishment of benchmark datasets such as MOT15, MOT16, and MOT17 has been instrumental in advancing the field, providing a framework for evaluating tracking methods and facilitating comparison between different approaches \cite{CVPR19 Tracking and Detection Challenge: How crowded can it get?}, \cite{MOT17: A Benchmark for Multi-Object Tracking}. The use of metrics such as precision, recall, and F1-score is common in evaluating tracking performance, although authors also discuss the limitations of these metrics and the need for more comprehensive evaluation frameworks \cite{Multiple Object Tracking: A Review}.

The development of standardized benchmarks has been complemented by the exploration of new tracking techniques, including deep learning approaches such as attention mechanisms, graph convolutional neural networks, and siamese networks \cite{Recent Advances in Multiple Object Tracking: Techniques and Applications}. These techniques have shown promise in addressing challenges such as occlusion, similar appearance, and ID switching, which are consistently highlighted as areas where current methods struggle \cite{Multiple Object Tracking: A Review}, \cite{Recent Advances in Multiple Object Tracking: Techniques and Applications}. The importance of considering factors such as object visibility and class labels is also emphasized in the development of benchmark datasets, highlighting the need for a nuanced understanding of the complexities involved in MOT \cite{MOTChallenge: A Benchmark for Multi-Object Tracking}.

Despite these advances, the papers acknowledge that current tracking methods still face significant challenges, including the need for more robust and efficient algorithms, particularly in applications such as autonomous driving \cite{Multiple Object Tracking: A Review}, \cite{Recent Advances in Multiple Object Tracking: Techniques and Applications}. The development of standardized benchmarks and evaluation methodologies is seen as an ongoing challenge, requiring continued effort to establish clear and effective metrics that can accurately capture the performance of MOT systems \cite{MOTChallenge: A Benchmark for Multi-Object Tracking}, \cite{CVPR19 Tracking and Detection Challenge: How crowded can it get?}. Furthermore, the papers highlight the importance of considering real-world scenarios and applications in the development of MOT systems, rather than solely focusing on benchmark performance \cite{Multiple Object Tracking: A Review}, \cite{Recent Advances in Multiple Object Tracking: Techniques and Applications}.

The use of leaderboards and ranking tables is seen as a way to provide a clear comparison of tracking methods, but authors also caution against over-emphasizing their importance, highlighting the need for a more nuanced understanding of the strengths and limitations of different approaches \cite{MOTChallenge: A Benchmark for Multi-Object Tracking}, \cite{CVPR19 Tracking and Detection Challenge: How crowded can it get?}. The potential applications of MOT, including autonomous driving and robot navigation, are frequently mentioned as motivations for advancing the field, emphasizing the need for continued innovation and development in this area \cite{Multiple Object Tracking: A Review}, \cite{Recent Advances in Multiple Object Tracking: Techniques and Applications}. Overall, the evaluation metrics and methodologies for MOT have undergone significant development in recent years, driven by the need for standardized benchmarks and objective measures of performance, and highlighting the importance of continued innovation and development in this field.

\section{Applications of MOT in Automated Driving Systems}
The applications of Multi-Object Tracking (MOT) in automated driving systems have garnered significant attention in recent years, with a plethora of research focusing on enhancing the accuracy, robustness, and efficiency of MOT algorithms in complex scenarios \cite{Multiple Object Tracking in Recent Times: A Literature Review}. A key challenge in this domain is the ability to effectively track multiple objects, including pedestrians, vehicles, and other obstacles, in real-time, while navigating through dynamic environments \cite{Multi-Object Tracking with Camera-LiDAR Fusion for Autonomous Driving}. To address these challenges, researchers have proposed various innovative approaches, including the fusion of camera and LiDAR data, which has been shown to improve tracking performance by leveraging the complementary strengths of different sensor modalities \cite{Know Your Surroundings: Panoramic Multi-Object Tracking by Multimodality Collaboration}.

The use of deep learning techniques has also become increasingly prevalent in MOT for automated driving, with many recent approaches leveraging transformer models, graph convolutional neural networks, and siamese networks to improve tracking accuracy by better handling appearance similarities and motion predictions \cite{MOTChallenge: A Benchmark for Single-Camera Multiple Target Tracking}. Furthermore, the importance of standardized evaluation frameworks has been emphasized, with benchmarks like MOTChallenge providing a comprehensive platform for assessing and comparing the performance of different MOT algorithms \cite{MOTChallenge: A Benchmark for Single-Camera Multiple Target Tracking}. These developments underscore the need for continued innovation in MOT, particularly in addressing challenges such as occlusion, similar object appearances, and varying lighting conditions, which are commonly encountered in real-world scenarios \cite{Multiple Object Tracking: A Literature Review}.

A notable trend in MOT research is the shift towards multimodal sensing, which combines data from different sensors to enhance robustness and accuracy \cite{Multi-Object Tracking with Camera-LiDAR Fusion for Autonomous Driving}. This approach has been shown to be particularly effective in complex environments, where single-modal sensing may struggle to provide reliable tracking performance. Moreover, the use of panoramic multi-object tracking frameworks, such as MMPAT, which leverages both 2D panorama images and 3D point clouds, has been proposed to enhance tracking performance in scenarios with significant occlusion or object density \cite{Know Your Surroundings: Panoramic Multi-Object Tracking by Multimodality Collaboration}. These advancements have significant implications for the development of automated driving systems, where reliable and accurate MOT is critical for ensuring safety and efficiency.

Despite these developments, several challenges persist, including the need for real-time processing, scalability, and generalizability across different scenarios \cite{Multiple Object Tracking in Recent Times: A Literature Review}. The computational efficiency of many advanced MOT algorithms remains a significant concern, particularly in resource-constrained environments. Additionally, the availability and quality of training data can significantly impact the performance of deep learning-based MOT approaches, highlighting the need for large, high-quality datasets and effective data augmentation strategies \cite{MOTChallenge: A Benchmark for Single-Camera Multiple Target Tracking}. Addressing these challenges will be crucial for realizing the full potential of MOT in automated driving systems and enabling the widespread adoption of autonomous vehicles.

In conclusion, the applications of MOT in automated driving systems have witnessed significant advancements in recent years, driven by innovations in multimodal sensing, deep learning techniques, and standardized evaluation frameworks. While challenges persist, the ongoing research efforts are expected to yield further improvements in the accuracy, robustness, and efficiency of MOT algorithms, ultimately contributing to the development of safer and more efficient automated driving systems \cite{Multiple Object Tracking: A Literature Review}. As the field continues to evolve, it is essential to address the remaining challenges and limitations, including scalability, real-time processing, and generalizability, to realize the full potential of MOT in automated driving systems.

\section{Advances in Detection-Aware and Polygon-Based Tracking}
The field of multi-camera multi-object tracking has witnessed significant advancements in recent years, with a particular focus on detection-aware and polygon-based tracking methods. These approaches have been designed to address the challenges and limitations of existing tracking algorithms, which often struggle to effectively handle complex scenarios such as background clutters and poor light conditions \cite{Challenge of Multi-Camera Tracking}. One notable development is the introduction of transformer-based approaches, such as MCTR, which leverages end-to-end detectors and probabilistic associations to enable robust multi-object detection and tracking across multiple cameras \cite{MCTR: Multi Camera Tracking Transformer}. This method has shown promising results, demonstrating the potential of deep learning-based techniques in enhancing tracking performance.

In addition to transformer-based methods, researchers have also explored the use of multimodality collaboration frameworks, such as MMPAT, which combines 2D panorama images and 3D point clouds to improve tracking accuracy in complex environments \cite{Know Your Surroundings: Panoramic Multi-Object Tracking by Multimodality Collaboration}. This approach highlights the importance of considering multiple modalities to enhance tracking performance, a theme that is echoed in several other papers \cite{Tracking the Trackers: An Analysis of the State of the Art in Multiple Object Tracking}. Furthermore, the use of bounding polygons for fast multi-object tracking and segmentation has been proposed in methods such as PolyTrack, which demonstrates the effectiveness of this approach in urban environment videos \cite{PolyTrack: Tracking with Bounding Polygons}.

A common theme across these papers is the need for more effective and efficient multi-camera tracking algorithms, which can handle complex scenarios and provide accurate tracking results. The use of end-to-end approaches, such as transformers and deep learning-based methods, is becoming increasingly popular in multi-object tracking research \cite{MCTR: Multi Camera Tracking Transformer}. However, these methods often require large amounts of labeled data for training, which can be time-consuming and expensive to obtain \cite{Challenge of Multi-Camera Tracking}. Moreover, the analysis of state-of-the-art trackers reveals current trends and weaknesses in multiple object tracking methods, providing guidance for future research \cite{Tracking the Trackers: An Analysis of the State of the Art in Multiple Object Tracking}.

The technical details of these methods vary, with MCTR using DETR as an end-to-end detector to produce detections and detection embeddings independently for each camera view \cite{MCTR: Multi Camera Tracking Transformer}. In contrast, PolyTrack detects objects by producing heatmaps of their center keypoint and computes a bounding polygon over each instance instead of traditional bounding boxes \cite{PolyTrack: Tracking with Bounding Polygons}. MMPAT's framework consists of four major modules: panorama image detection, multimodality data fusion, data association, and trajectory inference \cite{Know Your Surroundings: Panoramic Multi-Object Tracking by Multimodality Collaboration}. The evaluation metrics used to assess the performance of these methods also vary, with AP and MOTA being commonly used \cite{Tracking the Trackers: An Analysis of the State of the Art in Multiple Object Tracking}.

Despite the advancements in detection-aware and polygon-based tracking, there are still several limitations and challenges that need to be addressed. For example, handling complex scenarios such as background clutters and poor light conditions can affect tracking performance \cite{Challenge of Multi-Camera Tracking}. Additionally, the requirement for large amounts of labeled data for training can be a significant obstacle \cite{MCTR: Multi Camera Tracking Transformer}. Furthermore, the analysis of state-of-the-art trackers highlights the need for more standardized evaluation protocols and benchmarks to facilitate fair comparisons between different methods \cite{Tracking the Trackers: An Analysis of the State of the Art in Multiple Object Tracking}. Overall, the developments in detection-aware and polygon-based tracking demonstrate the ongoing efforts to improve multi-camera multi-object tracking, and further research is needed to address the remaining challenges and limitations in this field.

\section{Open Challenges and Future Directions in MOT Research}
The field of Multi-Object Tracking (MOT) has witnessed significant advancements in recent years, with various techniques and algorithms being proposed to improve tracking performance in diverse scenarios. Despite these efforts, several open challenges and future directions remain to be addressed in MOT research. A comprehensive analysis of relevant papers reveals key findings and contributions that highlight the importance of standardized benchmarks, such as MOTChallenge \cite{MOTChallenge}, in evaluating and comparing the performance of different MOT algorithms. The need for more robust and efficient tracking methods that can handle challenging scenarios, including crowded scenes, occlusions, and abrupt appearance changes, is also emphasized \cite{DeepLearningforMOT}.

The potential of using advanced techniques, such as attention mechanisms \cite{AttentionMechanisms}, graph convolutional neural networks \cite{GraphConvolutionalNeuralNetworks}, and siamese networks \cite{SiameseNetworks}, to improve MOT performance is also highlighted. Furthermore, the significance of evaluating MOT algorithms on diverse datasets and scenarios to ensure their generalizability and robustness cannot be overstated \cite{EvaluationMetricsforMOT}. The emphasis on benchmarking, focus on challenging scenarios, use of deep learning techniques, and importance of evaluation metrics are common themes that emerge from the analyzed papers.

Contrasting viewpoints or approaches are also present in the literature, with some papers comparing traditional MOT methods with deep learning-based approaches \cite{TraditionalvsDeepLearning}, while others discuss the differences between single-camera and multi-camera tracking \cite{SingleCameravsMultiCameraTracking}. The technical details and methodologies employed in the analyzed papers include the use of the MOTChallenge benchmark, various deep learning architectures, and different tracking algorithms, such as Kalman filter and particle filter \cite{MOTAlgorithms}. However, limitations and challenges remain, including handling crowded scenes, robustness to appearance changes, and scalability and efficiency \cite{LimitationsandChallengesinMOT}.

These findings and themes have significant implications for future research directions in MOT. The development of more efficient and scalable MOT algorithms that can handle large datasets and real-time applications is crucial \cite{EfficientandScalableMOT}. Moreover, the need to address challenging scenarios, such as crowded scenes and abrupt appearance changes, requires further investigation into advanced techniques and evaluation metrics \cite{AdvancedTechniquesforMOT}. The connections between MOT and other fields, such as computer vision and machine learning, also highlight the potential for interdisciplinary research and collaboration \cite{InterdisciplinaryResearchinMOT}. Ultimately, addressing these open challenges and future directions will be essential to improving the robustness, efficiency, and accuracy of MOT algorithms, with significant applications in areas such as surveillance, robotics, and autonomous vehicles. 

The references cited in this section are based on a comprehensive analysis of relevant papers \cite{Paper1, Paper2, Paper3, Paper4, Paper5, Paper6, Paper7, Paper8, Paper9, Paper10, Paper11, Paper12, Paper13, Paper14, Paper15, Paper16, Paper17, Paper18, Paper19, Paper20}, which demonstrate the significance of ongoing research in MOT and highlight areas where further investigation is needed to advance the field.

\section{Conclusion and Summary of Key Findings in MOT}
The survey on Multi Camera Multi Object Tracking has yielded significant insights into the current state of research in this field, with numerous papers contributing to the development of more effective tracking methods \cite{Multiple Object Tracking: A Literature Review}. The key findings from these papers highlight the importance of benchmarking efforts, such as MOTChallenge \cite{MOTChallenge: A Benchmark for Single-Camera Multiple Target Tracking}, which has established a standardized framework for evaluating Multi-Object Tracking (MOT) methods and has promoted progress in the field. Furthermore, the identification of challenges, including occlusions, appearance changes, and ID switching, emphasizes the need for robust tracking methods \cite{Challenge of Multi-Camera Tracking}.

A common theme throughout the analyzed papers is the focus on people tracking, particularly pedestrian tracking, due to its significance in applications such as autonomous driving and robot navigation \cite{MOT16: A Benchmark for Multi-Object Tracking}. The importance of developing tracking methods that can handle challenges such as occlusions and appearance changes is also widely recognized. In addition, the need for robust and efficient tracking algorithms is highlighted, with various approaches being discussed, including traditional methods, such as the Kalman filter and particle filter, and deep learning-based methods, including convolutional neural networks (CNNs) and recurrent neural networks (RNNs) \cite{Multiple Object Tracking in Recent Times: A Literature Review}.

The papers also present different viewpoints and approaches, with some focusing on single-camera tracking \cite{MOTChallenge: A Benchmark for Single-Camera Multiple Target Tracking}, while others discuss the unique challenges of multi-camera systems \cite{Challenge of Multi-Camera Tracking}. The technical details and methodologies employed in the papers are diverse, including evaluation metrics, such as MOTA and MOTP, to assess tracking performance \cite{MOT16: A Benchmark for Multi-Object Tracking}. However, limitations and challenges in MOT research are also highlighted, including the difficulties posed by occlusions and appearance changes, scalability and computational efficiency, and the lack of standardized datasets \cite{Multiple Object Tracking: A Literature Review}.

The implications of these findings are significant, as they highlight the need for continued research into more effective and efficient tracking methods. The connections between the different papers and approaches are also evident, with many building upon or extending existing work. For example, the development of deep learning-based methods has been influenced by traditional approaches, such as the Kalman filter \cite{Deep Learning for Multi-Object Tracking}. Overall, the survey provides a comprehensive overview of the current state of research in Multi Camera Multi Object Tracking, highlighting key themes, developments, and implications for future work. The analysis of the papers demonstrates that while significant progress has been made, there is still a need for more robust, efficient, and scalable tracking methods to address the challenges posed by real-world applications \cite{MOTChallenge: A Benchmark for Single-Camera Multiple Target Tracking}.

\section{Conclusion}
This survey on multi-camera multi-object tracking has provided an in-depth analysis of the current state of research in this field, highlighting key findings, common themes, and technical details from relevant papers \cite{Challenge of Multi-Camera Tracking}, \cite{Tracking the Trackers: An Analysis of the State of the Art in Multiple Object Tracking}, \cite{An equalised global graphical model-based approach for multi-camera object tracking}, \cite{MCTR: Multi Camera Tracking Transformer}, and \cite{MTMMC: A Large-Scale Real-World Multi-Modal Camera Tracking Benchmark}. The challenge of multi-camera tracking is distinct from single-camera tracking, with new technological and system architecture challenges that must be addressed. Standardized benchmarks, such as those proposed in \cite{Tracking the Trackers: An Analysis of the State of the Art in Multiple Object Tracking}, are crucial for evaluating multiple object tracking methods and have shown current trends and weaknesses in people tracking methods.

Global graph models, like the one presented in \cite{An equalised global graphical model-based approach for multi-camera object tracking}, can improve overall multi-camera object tracking performance by treating single camera tracking and inter-camera tracking differently. End-to-end approaches, such as the Multi-Camera Tracking Transformer (MCTR) \cite{MCTR: Multi Camera Tracking Transformer}, have also been shown to effectively detect and track objects across multiple cameras with overlapping fields of view. The importance of large-scale real-world datasets, like MTMMC \cite{MTMMC: A Large-Scale Real-World Multi-Modal Camera Tracking Benchmark}, cannot be overstated, as they provide a challenging test-bed for studying multi-camera tracking under diverse real-world complexities.

A common theme throughout the surveyed papers is the need for standardized benchmarks and evaluation frameworks to compare different multi-object tracking methods. Additionally, there is a growing interest in end-to-end approaches for multi-camera tracking, leveraging advances in deep learning and computer vision \cite{MCTR: Multi Camera Tracking Transformer}. The recognition of the challenges posed by real-world complexities, such as varying lighting conditions, weather, and seasons, is also a recurring theme \cite{MTMMC: A Large-Scale Real-World Multi-Modal Camera Tracking Benchmark}.

The technical details of the surveyed papers highlight the importance of considering both single camera tracking and inter-camera tracking in multi-camera tracking systems. Global graph models with improved similarity metrics can enhance multi-camera object tracking performance \cite{An equalised global graphical model-based approach for multi-camera object tracking}. End-to-end detectors like DEtector TRansformer (DETR) can be used to produce detections and detection embeddings independently for each camera view \cite{MCTR: Multi Camera Tracking Transformer}.

Despite the advances in multi-camera multi-object tracking, there are still limitations and challenges that must be addressed. The difficulty and cost of collecting and labeling large-scale real-world datasets for multi-camera tracking is a significant challenge \cite{MTMMC: A Large-Scale Real-World Multi-Modal Camera Tracking Benchmark}. Furthermore, the challenges posed by varying lighting conditions, weather, and seasons in real-world environments require more robust and efficient algorithms to handle the complexities of multi-camera tracking. As research continues to advance in this field, it is essential to consider these limitations and challenges to develop more effective and efficient multi-camera multi-object tracking systems. Overall, this survey has highlighted the key themes and developments in multi-camera multi-object tracking, discussing implications and connections between different approaches and techniques.

\end{document}